% !PRINTED 142
\weberpage{142}
%<<CONT

und demnach erhält man aus (6) das folgende von \gesperrt{Newton} herrührende Formelsystem\footnote[1]{\gesperrt{Newton}, Arithmetica universalis, edit. 's Gravesande, p.~592.}:
\begin{equation}
\begin{aligned}
0&=s_1+a_1,\\
0&=s_2+a_1s_1+2a_2,\\
0&=s_3+a_1s_2+a_2s_1+3a_3,\\
&\ \vdots\\
0&=s_{n-1}+a_1s_{n-2}+a_2s_{n-3}+\cdots+(n-1)a_{n-1}.
\end{aligned}
\tag{7}
\end{equation}

Dies Formelsystem lässt sich aber noch weiter fortsetzen; denn da
\[
 S[f(\alpha)]=0,\qquad S[\alpha f(\alpha)]=0,\qquad S[\alpha^2 f(\alpha)]=0,\ldots
\]
ist, so folgt
\begin{equation}
\begin{aligned}
0&=s_n+a_1s_{n-1}+a_2s_{n-2}+\cdots+n a_n,\\
0&=s_{n+1}+a_1s_n+a_2s_{n-1}+\cdots+a_ns_1,\\
0&=s_{n+2}+a_1s_{n+1}+a_2s_n+\cdots+a_ns_2,\quad\ldots
\end{aligned}
\tag{8}
\end{equation}

Durch die Formeln (7), (8) ist nun die Aufgabe gelöst, der Reihe nach die Functionen $s_1,s_2,s_3,\ldots$ bis zu beliebiger Höhe als ganze rationale Functionen der symmetrischen Grundfunctionen darzustellen, z. B.
\begin{equation}
\begin{aligned}
 s_1&=-a_1,\\
 s_2&=+a_1^2-2a_2,\\
 s_3&=-a_1^3+3a_1a_2-3a_3,\\
 s_4&=+a_1^4-4a_1^2a_2+4a_1a_3+2a_2^2-4a_4,\\
 &\hspace{1.5em}\vdots
\end{aligned}
\tag{9}
\end{equation}
und die Bildungsweise dieser Ausdrücke zeigt, dass die Coëfficienten in diesen Darstellungen \gesperrt{ganze Zahlen} sind.

Man kann auch umgekehrt mittelst der Formeln (7) und (8) die symmetrischen Grundfunctionen $a_1,a_2,a_3,\ldots$ rational durch die Potenzsummen $s_1,s_2,s_3,\ldots$ ausdrücken. Diese Darstellung ist aber insofern weit weniger einfach, als die Coëfficienten nicht ganze, sondern gebrochene Zahlen sind, z. B.
\begin{equation}
\begin{aligned}
 a_1&=-s_1,\\
 2a_2&=+s_1^2-s_2,\\
 6a_3&=-s_1^3+3s_1s_2-2s_3,\\
 &\hspace{1.5em}\vdots
\end{aligned}
\tag{10}
\end{equation}
